\section{Benchmark Scenario Specifications}
\label{app:scenarios}


This section provides detailed descriptions for the 7 benchmark scenarios summarized in Table~\ref{tab:scenarios_compact}. For each scenario, we describe the system objective, the constituent entities, and the governing dynamics.
We abopted this originally from \url{https://github.com/SimulationEverywhere/Cadmium-Simulation-Environment}, a widely acknowledged repository consisting various simulation tasks designed and implemented by human experts.

\subsection{Online Banking System (IOBS)}
\textbf{Domain:} Service Operations 
\textbf{Dynamics Type:} Pipeline with Probabilistic Routing

\begin{itemize}[leftmargin=*]
\item \textbf{Overview:} This scenario models a linear transaction processing pipeline for an online banking system. It simulate the sequential validation of user requests through multiple security and processing layers. 
\item \textbf{Entities:} The system consists of a chain of five modules: Account Access Manager, Account Number Verifier, Password Verifier, Bill Payment Manager, and Transaction Process Manager.
\item \textbf{Dynamics:} The system operates as a single-direction pipeline. A request enters at given time and passes through each stage with a fixed processing delay (10s).
\begin{itemize}
\item \textbf{Probabilistic Branching:} At the ANV and PV stages, requests have a 50$\%$ probability of failing verification. Failed requests are terminated immediately (dropped from the pipeline), while valid requests proceed.
\item \textbf{State Updates:} The final TPM stage maintains a persistent state (account balance) that is updated only upon the successful completion of the entire pipeline chain.
\end{itemize}
\end{itemize}

\subsection{O-Train Light Rail (OTrain)}
\textbf{Domain:} Transportation 
\textbf{Dynamics Type:} Schedule-driven with Serial Delays

\begin{itemize}[leftmargin=*]
\item \textbf{Overview:} This scenario simulates a light rail transit system with a single train shuttling bi-directionally between 5 stations (Bayview to Greenboro). It models the interaction between the train schedule and stochastic passenger arrival.
\item \textbf{Entities:} \textit{Train} (mobile resource), \textit{Stations} (fixed nodes), and \textit{Passengers} (transient entities).
\item \textbf{Dynamics:}
\begin{itemize}
\item \textbf{Schedule-Driven Movement:} The train adheres to a strict timetable, moving between stations with fixed travel intervals.
\item \textbf{Serial Serialization:} Unlike simple fluid queues, boarding and alighting are modeled as serial processes. Each passenger takes 0.025s to board or alight, creating a delay that is linearly proportional to the number of passengers.
\item \textbf{Stochastic Generation:} Passengers are generated at stations using a randomized interval (Normal distribution), with specific origin-destination constraints.
\end{itemize}
\end{itemize}

\subsection{Epidemic Compartmental Model (SEIRD)}
\textbf{Domain:} Biological Systems 
\textbf{Dynamics Type:} Continuous (ODE)

\begin{itemize}[leftmargin=*]
\item \textbf{Overview:} This scenario implements a classic SEIRD (Susceptible-Exposed-Infective-Recovered-Deceased) model to simulate disease spread within a closed population.
\item \textbf{Entities:} The population is treated as an aggregate divided into five compartments (S, E, I, R, D).
\item \textbf{Dynamics:} Unlike other event-driven scenarios, this model approximates continuous system dynamics using numerical integration (Forward Euler method).
\begin{itemize}
\item \textbf{State Evolution:} State updates occur at fixed time steps ( days). The transitions between compartments are governed by differential equations based on parameters such as transmission rate (), incubation period, and mortality rate.
\item \textbf{Conservation:} The model must strictly enforce population conservation at every time step despite floating-point operations.
\end{itemize}
\end{itemize}

\subsection{Offline File Transfer (OfflineFileTransfer)}
\textbf{Domain:} Network Systems 
\textbf{Dynamics Type:} Dual-Loop FSM with Interrupts

\begin{itemize}[leftmargin=*]
\item \textbf{Overview:} This scenario simulates a "Dropbox-like" synchronization system where a file upload process and a download process are decoupled by an intermediate server buffer.
\item \textbf{Entities:} \textit{Sender}, \textit{Server} (acting as a buffer), \textit{Receiver}, and four reliable \textit{Subnets} (channels).
\item \textbf{Dynamics:} The system features two asynchronous loops:
\begin{itemize}
\item \textbf{Dual-Loop Coupling:} The Upload Loop (Sender Server) and Download Loop (Server Receiver) operate independently but are coupled via the Server's storage queue. Both loops utilize the Alternating Bit Protocol (ABP) for reliability.
\item \textbf{External Interrupts:} The simulation is driven by a command stream. The request command acts as an external interrupt (or valve) that can dynamically pause or resume the download loop, requiring the system to handle state persistence during interruptions.
\end{itemize}
\end{itemize}

\subsection{C.5. Alternating Bit Protocol (ABP)}
\textbf{Domain:} Network Protocols 
\textbf{Dynamics Type:} Stop-and-Wait with Deterministic Noise

\begin{itemize}[leftmargin=*]
\item \textbf{Overview:} A standard implementation of the Alternating Bit Protocol (ABP) to achieve reliable data transmission over an unreliable channel.
\item \textbf{Entities:} \textit{Sender}, \textit{Receiver}, and two \textit{Subnets} (Forward and Backward channels).
\item \textbf{Dynamics:}
\begin{itemize}
\item \textbf{State Machine:} The Sender implements a "Stop-and-Wait" FSM. It sends a packet and blocks until a valid ACK (with the matching bit) is received or a timeout occurs.
\item \textbf{Deterministic Noise:} Packet loss is not random but governed by a Linear Congruential Generator (LCG) inside the subnets. This ensures the simulation is strictly reproducible: packet fate (drop vs. pass) is pre-determined by the seed, testing the model's ability to implement precise algorithmic logic.
\end{itemize}
\end{itemize}

\subsection{C.6. Strategic Airlift Operations (StrategicAirlift)}
\textbf{Domain:} Logistics 
\textbf{Dynamics Type:} Active Reneging \& Resource Cycling

\begin{itemize}[leftmargin=*]
\item \textbf{Overview:} A complex logistics simulation managing the transport of cargo pallets via a fleet of aircraft.
\item \textbf{Entities:} \textit{Facility} (Source), \textit{Loading Queue}, \textit{Fleet Coordinator}, multiple \textit{Aircraft}, and \textit{Destination}.
\item \textbf{Dynamics:}
\begin{itemize}
\item \textbf{Active Reneging:} Pallets in the queue have a finite lifespan. The queue must actively monitor expiration times and autonomously discard (renege) items that stay too long, even without external events.
\item \textbf{Resource Scheduling:} Aircraft follow a multi-stage state cycl. The Coordinator must match available cargo with idle aircraft, handling resource contention.
\end{itemize}
\end{itemize}

\subsection{C.7. Barbershop (Barbershop)}
\textbf{Domain:} Service Queueing 
\textbf{Dynamics Type:} Blocking \& Handshake Signaling

\begin{itemize}[leftmargin=*]
\item \textbf{Overview:} A classic queuing simulation of a barbershop with a reception area and a two-stage service process (Inspection and Cutting).
\item \textbf{Entities:} \textit{Reception Desk}, \textit{Inspection Station}, and \textit{Cutting Station}.
\item \textbf{Dynamics:}
\begin{itemize}
\item \textbf{Finite Blocking:} The reception queue has a hard capacity (limit 8). Arrivals when the queue is full are strictly rejected (Blocking).
\item \textbf{Inter-module Signaling:} The two service stages (Inspection and Cutting) operate sequentially. The Inspection module must wait for a "done" signal from the Cutting module (Handshake) before it can release the current customer and accept a new one from Reception.
\end{itemize}
\end{itemize}